\section{Introduction}\label{sec:intro}

Sparse matrix--vector multiplication (SpMV, $y = Ax$) and sparse matrix--matrix multiplication (SpGEMM,
$C = AB$) dominate iterative solvers, multigrid set-up and graph analytics. Both are limited by memory traffic
rather than arithmetic, and both are irregular: how evenly the nonzeros are spread over rows, how long the longest
row is, and how large the operands are relative to GPU memory all change which implementation is fastest.

The two papers this project starts from show this clearly, and both stop short of acting on it.
Heller and Oberhuber's adaptive row-grouped CSR (ArgCSR) gives long rows more threads, and its performance depends
strongly on one parameter, \texttt{desiredChunkSize}: on \texttt{rajat23} a value of~1 is six times faster than~32,
while the AFE matrices prefer~32. The paper then fixes the value to~1 for every matrix, and concludes that ``if high
performance is the top priority, one should test more formats and choose the best one''~\cite{heller2012argcsr}.
Deveci et al.\ extend the KokkosKernels SpGEMM to multilevel memory: when the operands do not fit in
fast memory they keep one operand's row block resident and stream the other through it~\cite{deveci2018multilevel}.
Their partitioning heuristic (Algorithm~4) reserves a fixed 75\%/25\% of fast memory for the resident and the
streamed operand and chooses between the two chunking variants by copy volume alone. Copies and computation are not
overlapped, which the authors leave as future work.

RASP (\emph{Runtime-Adaptive Sparse Processing}) makes these choices at run time, per matrix, from cheap
measurements:
\begin{enumerate}
\item \textbf{A lightweight SpMV selector.} One GPU kernel extracts structural features; a shallow decision tree,
  trained offline and compiled into the library as plain C++, predicts the best format \emph{and} its parameters;
  a break-even test converts formats only when the expected number of calls repays the conversion (\cref{sec:design-spmv}).
\item \textbf{An overlapped, self-calibrating out-of-core SpGEMM.} Paper~1's chunking, extended with double and
  triple buffering, a cost model that simulates the copy/compute pipeline using bandwidths and kernel rates
  measured at run time, and re-partitioning while it runs (\cref{sec:design-spgemm}).
\item \textbf{A re-evaluation of both papers} on a PCIe-attached consumer GPU (RTX~4050 Laptop, Ada) under
  WSL2, against the 2026 cuSPARSE (\cref{sec:eval}).
\end{enumerate}
Automatic format selection with machine learning is an established research area
(\cref{sec:related}); our novelty is relative to the two assigned papers: tuning ArgCSR's parameters per matrix and
per row region, a conversion-aware cost model, and bringing runtime adaptivity and overlap to out-of-core SpGEMM.

\Cref{tab:goals} lists the goals we set at the start and where each is evaluated.

\begin{table}[t]
\centering\small
\caption{Project goals and where they are evaluated.}\label{tab:goals}
\begin{tabular}{@{}lp{6.2cm}p{5.6cm}l@{}}
\toprule
 & Goal & Target & Section \\
\midrule
\G1 & Reproduce paper 2 (ArgCSR) on modern hardware & same qualitative winners and losers & \ref{sec:eval-g1} \\
\G2 & Runtime SpMV selector & geomean $T_\text{selected}/T_\text{oracle} \le 1.10$; never $>$5\% slower than cuSPARSE on matrices $\ge$ 1M nnz & \ref{sec:eval-g2} \\
\G3 & Selection is cheap & features + inference $\le$ 2 SpMVs & \ref{sec:eval-g2} \\
\G4 & Conversion-aware choice & never $>$5\% worse than cuSPARSE at $N=1$ & \ref{sec:eval-g2} \\
\G5 & Paper 1's placement study on a PCIe GPU & our version of its Table 3 & \ref{sec:eval-g5} \\
\G6 & Overlapped, adaptive out-of-core SpGEMM & $\ge 1.3\times$ over Algorithm 4 when the problem exceeds the budget & \ref{sec:eval-g6} \\
\G7 & End-to-end applications & CG (SpMV) and AMG set-up (SpGEMM) working and timed & \ref{sec:eval-g7} \\
\bottomrule
\end{tabular}
\end{table}
